%% input/0_introduction.tex
\section{Introduction}
\label{sec:introduction}

% Introduce the problem, motivation, and contributions here.

    The recent transformation within the field of High Performance Computing (HPC) towards heterogeneous architectures has 
    introduced issues such as the curse of dimensionality and the exponential scaling of Hilbert space, 
    causing the focus to shift towards the potential advantages of integrating Quantum Processing Units (QPUs). 
    QPUs have emerged, not as the replacement for classical computation, but 
    instead as specialized accelerators to be integrated into our existing classical systems, capable of handling 
    specific, formerly intractible, computational tasks ~\parencite{rallis2025interfacingquantumcomputingsystems}. This has lead to the 
    development of hybrid classical-quantum frameworks so to achieve new grounds in the fields of material science, 
    molecular simulation, complex optimization, Hamiltonian simulation, alongside many others.   

%     --> MENTION HOW WE REMOVE THE OPERATIONAL OVERHEAD FROM RESEARCHERS AND QPU/DATACENTER DEVELOPERS BY ORCHESTRATING THE COORDINATION OF CLOUD HARDWARES
%     --> Hamiltonian simulation, optimization, and quantum machine learning, with broader
% application areas including electronic structure and related chemistry, materials, and data-analysis
% workloads  (https://arxiv.org/pdf/2604.20912)

    In our current NISQ era we have at our disposal quantum hardware that is limited and noisy, not yet 
    able to achieve the full effects of the "quantum advantage" ~\parencite{harrow2017quantumsupremacy}. The most promising path 
    is hybrid algorithms that combine each system's computational benefits,   
    specifically the Variational Quantum Eigensolver (VQE), which continually offloads heavy computation to its classical resources while 
    QPUs handle circuit execution ~\parencite{peruzzo2014variational, rallis2025interfacingquantumcomputingsystems} so to bypass their hardware limitations. 
    Alone, classical HPC structures struggle when the problem dimension grows exponentially (PLACE THE NUMERICAL REFERENCE HERE) alongside 
    the size of the physical system being observed, prompting the exploration of quantum hardware as accelerators. is handled by QPUs 

    Despite the theoretical potential of hybrid collaboration, there is a significant architectural gap, as  
    there is no specialized middleware to bridge remote quantum and classical components alongside distributed and accelerated HPC resources. Classical HPC frameworks (MPI, CUDA)
    are not designed to communicate directly with quantum control stacks and other existing quantum frameworks (Qiskit, Cirq, EXTEND REFERENCES HERE - IS IT BAD FAITH TO INCLUDE XACC?), which  
    lack features to effectively integrate with distributed HPC environments. Without a standardized hybrid middleware, researchers face issues of high tail latencies when synchronizing
    HPC resources and cloud-based QPUs, leading to performance bottlenecks that negate the possibility of a computational speedup offered by a
    hybrid system. 

    -> INTRODUCE AND MENTION THE COMPARISONS HERE BETWEEN PENNYLANE LIGHTNING, QISKIT AER MPI


    This paper presents a hybrid quantum-classical software stack that builds the middleware layer needed
    to abstract hardware and automate the dispatching of tasks. Existing middlewares, such as XACC, operate synchronously within a single host, whereas this stack introduces MPI distributed parallelization alongside 
    GPU acceleration, asynchronous QPU dispatching, and an intuitive API that abstracts the cloud-based QPU backend. The stack is benchmarked against VQE workflows 
    on a single host Docker environment with MPI/GPU acceleration aiming to demonstrate measurable middleware speedup, shown through the combined GPU and MPI path reaching of a (2-2.5X SPEEDUP NOW, NO?)x 
    speedup in the end-to-end wall-clock time when being compared to the baselines of serial, non-distributed classical implementations. 

    --> SHOULD I MENTION MY SPECIFIC HARDWARES HERE, NOW, OR LATER IN THE PAPER SO TO STRESS THAT IT IS HARDWARE/CLOUD PLUGIN AGNOSTIC?
    --> I WILL ALSO BE DOING A COMPARISON BETWEEN PENNYLANE LIGHTNING, QISKIT AER MPI

    The outline of the following sections in this paper is as follows: Section 2
    provides the system features and architecture, experimental configuration, and evaluation methodology (including compairsons to other works.). Section 4 expands on the experimental procedures with their results presented within Section 5. 
    Section 6 is the discussion of these findings and their implications, Section 7 covers future work, and Section 8 provides the thesis conclusion. 